\subsection{Künneth formula}

In this number, we consider a complex (C, d) of right A-modules and a complex (C', d') of left A-modules. Consider the canonical exact sequences

\[
0 \rightarrow \mathrm{Z} (\mathrm{C}) \stackrel {{j}} {{\rightarrow}} \mathrm{C} \stackrel {{\delta}} {{\rightarrow}} \mathrm{B} (\mathrm{C}) (- 1) \rightarrow 0,\tag{I}
\]

\[
0 \rightarrow \mathrm{B} (\mathrm{C}) \xrightarrow {i} \mathrm{Z} (\mathrm{C}) \xrightarrow {p} \mathrm{H} (\mathrm{C}) \rightarrow 0;\tag{II}
\]

we deduce from δ a k-homomorphism

\[
\mathrm{H} (\delta \otimes 1): \mathrm{H} (\mathrm{C} \otimes_ {\mathrm{A}} \mathrm{C} ^ {\prime}) \rightarrow \mathrm{H} (\mathrm{B} (\mathrm{C}) \otimes_ {\mathrm{A}} \mathrm{C} ^ {\prime}) (- 1);
\]

we deduce from (II) a connecting homomorphism

\[
\partial (\mathrm{II}, \mathrm{H} (\mathrm{C} ^ {\prime})): \operatorname{Tor} _ {1} ^ {\mathbf {A}} (\mathrm{H} (\mathrm{C}), \mathrm{H} (\mathrm{C} ^ {\prime})) \rightarrow \mathrm{B} (\mathrm{C}) \otimes_ {\mathbf {A}} \mathrm{H} (\mathrm{C} ^ {\prime});
\]

if we equip \(\operatorname{Tor}_{1}^{\mathbf{A}}(\mathrm{H}(\mathbf{C}),\mathrm{H}(\mathbf{C}^{\prime}))\) with the grading whose homogeneous component of degree \(n\) is \(\bigoplus_{p+q=n}\operatorname{Tor}_{1}^{\mathbf{A}}(\mathrm{H}_{p}(\mathbf{C}),\mathrm{H}_{q}(\mathbf{C}^{\prime}))\) , this connecting homomorphism is graded of degree 0. Moreover, there is a canonical homomorphism (X, p. 62)

\[
\gamma (\mathrm{B} (\mathrm{C}), \mathrm{C} ^ {\prime}): \mathrm{B} (\mathrm{C}) \otimes_ {\mathrm{A}} \mathrm{H} (\mathrm{C} ^ {\prime}) \rightarrow \mathrm{H} (\mathrm{B} (\mathrm{C}) \otimes_ {\mathrm{A}} \mathrm{C} ^ {\prime}).
\]

With these notations:

THEOREM 3. --- Suppose the A-modules B(C) and Z(C) are flat. There exists a unique homomorphism of graded k-modules, of degree -1,

\[
\alpha : \mathrm{H} (\mathrm{C} \otimes_ {\mathrm{A}} \mathrm{C} ^ {\prime}) \rightarrow \operatorname{Tor} _ {1} ^ {\mathrm{A}} (\mathrm{H} (\mathrm{C}), \mathrm{H} (\mathrm{C} ^ {\prime}))
\]

thereby making the diagram commutative

\[
\begin{array}{c} \mathrm{H} (\mathrm{C} \otimes_ {\mathrm{A}} \mathrm{C} ^ {\prime}) \xrightarrow {\alpha} \operatorname{Tor} _ {1} ^ {\mathrm{A}} (\mathrm{H} (\mathrm{C}), \mathrm{H} (\mathrm{C} ^ {\prime})) (- 1) \\ \mathrm{H} (\delta \otimes 1) \Biggl \downarrow \qquad \qquad \qquad \qquad \qquad \qquad \qquad \qquad \qquad \qquad \qquad \qquad \Biggl \downarrow \hat {c} (\Pi , \mathrm{H} (\mathrm{C} ^ {\prime})) \\ \mathrm{H} (\mathrm{B} (\mathrm{C}) \otimes_ {\mathrm{A}} \mathrm{C} ^ {\prime}) (- 1) \xleftarrow {\gamma (\mathrm{B} (\mathrm{C}) . \mathrm{C} ^ {\prime})} (\mathrm{B} (\mathrm{C}) \otimes_ {\mathrm{A}} \mathrm{H} (\mathrm{C} ^ {\prime})) (- 1). \end{array}
\]

The sequence of graded k-modules

\[
0 \longrightarrow \mathrm{H} (\mathrm{C}) \otimes_ {\mathrm{A}} \mathrm{H} \left(\mathrm{C} ^ {\prime}\right) \xrightarrow {\gamma \left(\mathrm{C} , \mathrm{C} ^ {\prime}\right)} \mathrm{H} \left(\mathrm{C} ^ {\prime} \otimes_ {\mathrm{A}} \mathrm{C} ^ {\prime}\right) \xrightarrow {\alpha} \operatorname{Tor} _ {1} ^ {\mathrm{A}} \left(\mathrm{H} (\mathrm{C}), \mathrm{H} \left(\mathrm{C} ^ {\prime}\right)\right) (- 1) \longrightarrow 0
\]

is exact.

Thus for each \( n \), we have an exact sequence

\[
\begin{array}{c} 0 \longrightarrow \bigoplus_ {p + q = n} \mathrm{H} _ {p} (\mathrm{C}) \otimes_ {\mathrm{A}} \mathrm{H} _ {q} (\mathrm{C} ^ {\prime}) \xrightarrow {\gamma_ {n} (\mathrm{C} , \mathrm{C} ^ {\prime})} \mathrm{H} _ {n} (\mathrm{C} \otimes_ {\mathrm{A}} \mathrm{C} ^ {\prime}) \\ \xrightarrow {\alpha_ {n}} \bigoplus_ {p + q = n - 1} \operatorname{Tor} _ {1} ^ {\mathrm{A}} \left(\mathrm{H} _ {p} (\mathrm{C}), \mathrm{H} _ {q} (\mathrm{C} ^ {\prime})\right) \longrightarrow 0. \end{array}\tag{11}
\]

For simplicity set \(\mathbf{B} = \mathbf{B}(\mathbf{C})\) , \(\mathbf{Z} = \mathbf{Z}(\mathbf{C})\) , \(\mathbf{H} = \mathbf{H}(\mathbf{C})\) and \(\mathbf{H}' = \mathbf{H}(\mathbf{C}')\) :

Since B is flat, we deduce from (I) an exact sequence (X, p. 72, cor. 2)

\[
0 \longrightarrow Z \otimes_ {A} C ^ {\prime} \xrightarrow {j \otimes 1} C \otimes_ {A} C ^ {\prime} \xrightarrow {\delta \otimes 1} (B \otimes_ {A} C ^ {\prime}) (- 1) \longrightarrow 0.\tag{12}
\]

Lemma 3. --- The connecting homomorphism H(B ⊗\_A C') → H(Z ⊗\_A C') associated with the exact sequence (12) is equal to H(i ⊗ 1).

Indeed, let \(a \in \mathbf{Z}(\mathbf{B} \otimes_{\mathbf{A}} \mathbf{C}')\) ; since B is flat, \(a\) belongs to the image of \(\mathbf{B} \otimes_{\mathbf{A}} \mathbf{Z}(\mathbf{C}')\) , hence can be written as \(\sum_{\lambda} da_{\lambda} \otimes b_{\lambda}\) , with \(a_{\lambda} \in \mathbf{C}\) , \(b_{\lambda} \in \mathbf{C}'\) , \(db_{\lambda} = 0\) . The image of

the class of \(a\) under the sought homomorphism is by definition the class of \(\mathbf{D}(\sum a_{\lambda}\otimes b_{\lambda})=\sum da_{\lambda}\otimes b_{\lambda}=(i\otimes1)(a)\) , whence the lemma.

The exact homology sequence associated with (12) is therefore

\[
\begin{array}{r l} \mathrm{H} (\mathrm{B} \otimes_ {\mathrm{A}} \mathrm{C} ^ {\prime}) & \xrightarrow {\mathrm{H} (i \otimes 1)} \mathrm{H} (\mathrm{Z} \otimes_ {\mathrm{A}} \mathrm{C} ^ {\prime}) \xrightarrow {\mathrm{H} (j \otimes 1)} \mathrm{H} (\mathrm{C} \otimes_ {\mathrm{A}} \mathrm{C} ^ {\prime}) \\ & \xrightarrow {\mathrm{H} (\delta \otimes 1)} \mathrm{H} (\mathrm{B} \otimes_ {\mathrm{A}} \mathrm{C} ^ {\prime}) (- 1) \xrightarrow {\mathrm{H} (i \otimes 1)} \mathrm{H} (\mathrm{Z} \otimes_ {\mathrm{A}} \mathrm{C} ^ {\prime}) (- 1). \end{array}
\]

Moreover, since Z is flat, we deduce from (II) an exact sequence of graded k-modules

\[
0 \longrightarrow \operatorname{Tor} _ {1} ^ {\mathrm{A}} (\mathrm{H}, \mathrm{H} ^ {\prime}) \xrightarrow {\hat {c} (\mathrm{II} , \mathrm{H} ^ {\prime})} \mathrm{B} \otimes_ {\mathrm{A}} \mathrm{H} ^ {\prime} \xrightarrow {i \otimes 1} Z \otimes_ {\mathrm{A}} \mathrm{H} ^ {\prime} \xrightarrow {p \otimes 1} \mathrm{H} \otimes_ {\mathrm{A}} \mathrm{H} ^ {\prime} \longrightarrow 0;
\]

finally, we have the canonical homomorphisms of no. 1

\[
\gamma_ {\mathbf {B}} = \gamma (\mathbf {B}, \mathbf {C} ^ {\prime}): \mathbf {B} \otimes_ {\mathbf {A}} \mathrm{H} ^ {\prime} \rightarrow \mathrm{H} (\mathbf {B} \otimes_ {\mathbf {A}} \mathbf {C} ^ {\prime})
\]

\[
\gamma_ {Z} = \gamma (Z, C ^ {\prime}): Z \otimes_ {A} H ^ {\prime} \rightarrow H (Z \otimes_ {A} C ^ {\prime})
\]

\[
\gamma_ {\mathrm{C}} = \gamma (\mathrm{C}, \mathrm{C} ^ {\prime}): \mathrm{H} \otimes_ {\mathrm{A}} \mathrm{H} ^ {\prime} \rightarrow \mathrm{H} (\mathrm{C} \otimes_ {\mathrm{A}} \mathrm{C} ^ {\prime}),
\]

whence a diagram of graded k-modules, with exact rows

\[
\begin{array}{r l} & \mathbf {B} \otimes \mathbf {H} ^ {\prime} \xrightarrow [ \gamma_ {B} ]{i \otimes 1} \mathbf {Z} \otimes \mathbf {H} ^ {\prime} \xrightarrow [ \gamma_ {Z} ]{p \otimes 1} \mathbf {H} \otimes \mathbf {H} ^ {\prime} \longrightarrow 0 \\ & \mathrm{H} (\mathbf {B} \otimes \mathbf {C} ^ {\prime}) \xrightarrow {\mathrm{H} (i \otimes 1)} \mathrm{H} (\mathbf {Z} \otimes \mathbf {C} ^ {\prime}) \xrightarrow {\mathrm{H} (j \otimes 1)} \mathrm{H} (\mathbf {C} \otimes \mathbf {C} ^ {\prime}) \xrightarrow {\mathrm{H} (\delta \otimes 1)} \mathrm{H} (\mathbf {B} \otimes \mathbf {C} ^ {\prime}) (- 1) \xrightarrow {\mathrm{H} (i \otimes 1)} \mathrm{H} (\mathbf {Z} \otimes \mathbf {C} ^ {\prime}) (- 1) \\ & 0 \longrightarrow \operatorname{Tor} _ {1} ^ {\mathrm{A}} (\mathrm{H}, \mathrm{H} ^ {\prime}) (- 1) \xrightarrow {\delta (\mathrm{II.H} ^ {\prime})} (\mathbf {B} \otimes \mathbf {H} ^ {\prime}) (- 1) \xrightarrow [ i \otimes 1 ]{\gamma_ {B}} (\mathbf {Z} \otimes \mathbf {H} ^ {\prime}) (- 1), \end{array}
\]

which is commutative by definition of homomorphisms \(\gamma\). But, since the complexes B and Z are split and flat, \(\gamma_{\mathbf{B}}\) and \(\gamma_{\mathbf{Z}}\) are bijective (X, p. 65, cor. 1). From this one deduces, on the one hand, that \(\gamma_{\mathbf{C}}\) is injective with image equal to Ker H( \(\delta \otimes 1\) ), and on the other hand that \(\gamma_{\mathbf{B}} \circ \partial(\mathrm{II}, \mathrm{H}')\) is injective, with image equal to Im H( \(\delta \otimes 1\) ). The theorem follows immediately from this.

COROLLARY 1. --- If B(C) and Z(C) are flat, then for every left A-module N and every integer n there is an exact sequence

\[
0 \longrightarrow \mathrm{H} _ {n} (\mathrm{C}) \otimes_ {\mathrm{A}} \mathrm{N} \xrightarrow {\gamma_ {n}} \mathrm{H} _ {n} (\mathrm{C} \otimes_ {\mathrm{A}} \mathrm{N}) \xrightarrow {\alpha_ {n}} \operatorname{Tor} _ {1} ^ {\mathrm{A}} \left(\mathrm{H} _ {n - 1} (\mathrm{C}), \mathrm{N}\right) \longrightarrow 0.\tag{13}
\]

COROLLARY 2. --- Suppose B(C) and B(C') are projective and Z(C) is flat. Then the sequences of k-modules (11) and (13) are exact and split.

This follows from the theorem and the following lemma:

Lemma 4. --- If B(C) and B(C') are projective, then the canonical homomorphism

\[
\gamma (\mathbf {C}, \mathbf {C} ^ {\prime}): \mathrm{H} (\mathbf {C}) \otimes_ {\mathbf {A}} \mathrm{H} (\mathbf {C} ^ {\prime}) \rightarrow \mathrm{H} (\mathbf {C} \otimes_ {\mathbf {A}} \mathbf {C} ^ {\prime})
\]

has a k-linear retraction.

Indeed, according to X, p. 65, remark b), there exist homologisms \(\varphi: C \to H(C)\) and \(\varphi': C' \to H(C')\) such that \(H(\varphi) = 1_{H(C)}\) and \(H(\varphi') = 1_{H(C')}\) . In the commutative diagram

\[
\begin{array}{c} \mathrm{H(C)} \otimes_ {\mathrm{A}} \mathrm {H(C^ {\prime})} \xrightarrow {\gamma (\mathrm{C,C} ^ {\prime})} \mathrm {H(C} \otimes_ {\mathrm{A}} \mathrm {C^ {\prime})} \\ \mathrm {H}(\varphi) \otimes \mathrm {H}(\varphi^ {\prime}) \Big \downarrow \\ \mathrm {H(C)} \otimes_ {\mathrm{A}} \mathrm {H(C^ {\prime})} \xrightarrow {\gamma (\mathrm{H(C),H(C^ {\prime})})} \mathrm {H(C)} \otimes_ {\mathrm{A}} \mathrm {H(C^ {\prime})} \end{array}
\]

\(\mathrm{H}(\varphi) \otimes \mathrm{H}(\varphi') \text{ and } \gamma(\mathrm{H}(\mathrm{C}), \mathrm{H}(\mathrm{C}') ) \text{ are the identity, whence the assertion.}\)

COROLLARY 3 (“universal coefficient formula”). --- Suppose the ring A is principal. If the complexes C and C' are free, then the sequences of A-modules (11) are exact and split; if the complex C is free, then the sequences of A-modules (13) are exact and split for every A-module N.

Indeed, B(C), Z(C), and B(C') are submodules of the free modules C, C, C', hence are free (VII, § 3, cor. 2 to th. 1), and one applies cor. 2.

COROLLARY 4 (“Künneth formula”). --- Suppose C is bounded above, and C and H(C) are flat; then the canonical homomorphism

\[
\gamma (\mathrm{C}, \mathrm{C} ^ {\prime}): \mathrm{H} (\mathrm{C}) \otimes_ {\mathrm{A}} \mathrm{H} (\mathrm{C} ^ {\prime}) \rightarrow \mathrm{H} (\mathrm{C} \otimes_ {\mathrm{A}} \mathrm{C} ^ {\prime})
\]

is bijective.

By the theorem, it suffices to prove that B(C) and Z(C) are flat. Now we have exact sequences

\[
0 \to \mathrm{B} _ {n} (\mathrm{C}) \to \mathrm{Z} _ {n} (\mathrm{C}) \to \mathrm{H} _ {n} (\mathrm{C}) \to 0
\]

\[
0 \to Z _ {n} (\mathbf {C}) \to \mathbf {C} _ {n} \to \mathbf {B} _ {n - 1} (\mathbf {C}) \to 0,
\]

whence, by X, p. 75, cor. 1, the implications \((\mathbf{B}_{n-1}(\mathbf{C})\) is flat) \(\Rightarrow (\mathbf{Z}_n(\mathbf{C})\) is flat) \(\Rightarrow (\mathbf{B}_n(\mathbf{C})\) is flat); we conclude by noting that \(\mathbf{B}_n(\mathbf{C}) = 0\) for \(n\) small enough.

COROLLARY 5. --- Let \(u: \mathbf{C} \to \mathbf{C}'\) a homomorphism of complexes of right A-modules, flat and bounded on the right. For every complex E of left A-modules, the morphism \(u \otimes 1_{\mathbf{E}}: \mathbf{C} \otimes_{\mathbf{A}} \mathbf{E} \to \mathbf{C}' \otimes_{\mathbf{A}} \mathbf{E}\) is a homologism.

Indeed, Con \((u)\) is a flat complex, bounded on the right and with zero homology; we therefore have H(Con \((u) \otimes_{\mathbf{A}} \mathbf{E}) = 0\) by cor. 4, hence H(Con \((u \otimes 1_{\mathbf{E}})) = 0\) (X, p. 67, lemma 2), and \(u \otimes 1_{\mathbf{E}}\) is a homologism.

